%---------------------------------------------------------------------
\chapter{Agent Design and Architecture}
%---------------------------------------------------------------------

This chapter presents the architecture of LoomTrace, describes each component shown in the
architecture diagram, justifies the agentic patterns that were selected against the alternatives that
were rejected, and describes the control flow, the shared state and the flow of data through the
system for one concrete input.

\section{System Architecture}

Figure~\ref{fig:arch} shows the architecture of LoomTrace. A product claim enters at the top left and
passes through six nodes of a state graph: \emph{Extract}, \emph{Planner}, the \emph{Retriever agent},
the \emph{Verifier agent}, the \emph{Grader} and the \emph{Guardrail}. The retriever agent calls typed
tools that read four data stores. The grader decides whether the evidence is sufficient; if it is not,
control returns to the planner, at most three times in total. The guardrail issues the verdict with its
citations, or escalates the case to a review queue, from which a human reviewer records the final
decision. All nodes read and write a single shared state object, \texttt{CaseState}, which is saved to
an SQLite database after every node.

\begin{figure}[H]
\centering
\resizebox{\linewidth}{!}{%
\begin{tikzpicture}[
  font=\small, every node/.append style={execute at begin node=\setstretch{1.0}},
  io/.style={draw=gray!70,fill=gray!8,rounded corners,minimum width=2.3cm,minimum height=1.0cm,align=center},
  slm/.style={draw=Plum,fill=Plum!8,rounded corners,minimum width=2.3cm,minimum height=1.0cm,align=center},
  agent/.style={draw=bmsblue,fill=guidebg,rounded corners,line width=1pt,minimum width=2.3cm,minimum height=1.0cm,align=center},
  rule/.style={draw=ForestGreen,fill=dobg,minimum width=2.3cm,minimum height=1.0cm,align=center},
  tool/.style={draw=ForestGreen,fill=dobg,minimum width=5.6cm,minimum height=1.0cm,align=center},
  store/.style={draw=gray,fill=gray!10,cylinder,shape border rotate=90,aspect=0.18,minimum width=2.4cm,minimum height=1.1cm,align=center,font=\footnotesize},
  human/.style={draw=BrickRed,fill=dontbg,rounded corners,minimum width=2.3cm,minimum height=1.0cm,align=center},
  dec/.style={draw=bmsblue,fill=white,diamond,aspect=1.6,inner sep=1pt,align=center,font=\footnotesize},
  arr/.style={-{Stealth[length=2.2mm]},thick},
  lab/.style={font=\scriptsize,fill=white,inner sep=1pt}]

% data stores (top)
\node[store] (reg) at (9.0,4.8) {GI registry\\(IP India +\\PIB catalogue)};
\node[store] (docs) at (12.25,4.8) {GI passages\\(BM25 index)};
\node[store] (cen) at (15.5,4.8) {Census 2019--20;\\weaver + batch\\(simulated)};
% tools
\node[tool,minimum width=6.8cm] (tools) at (12.25,2.5) {Typed tools (Pydantic): \texttt{gi\_registry\_lookup},\\\texttt{gi\_doc\_search}, \texttt{census\_lookup}, \texttt{weaver\_lookup},\\\texttt{provenance\_lookup}};
% main row
\node[io] (in) at (0,0) {Product claim\\(label text /\\image)};
\node[slm] (ext) at (3.5,0) {Extract\\(parser / OCR\\+ SLM)};
\node[slm] (plan) at (7.0,0) {Planner\\(checks + query)};
\node[agent] (ret) at (10.5,0) {Retriever\\agent};
\node[agent] (ver) at (14.0,0) {Verifier\\agent};
% second row
\node[rule] (gr) at (14.0,-2.9) {Grader\\(grade 0--1)};
\node[dec] (dec) at (10.0,-2.9) {Enough\\evidence?};
\node[rule] (gd) at (5.5,-2.9) {Guardrail\\+ HITL gate};
\node[io] (out) at (0,-2.9) {Cited verdict\\(3 classes)};
% third row
\node[store] (q) at (5.5,-5.6) {Review queue\\(SQLite)};
\node[human] (hum) at (0,-5.6) {Human\\reviewer};
\node[store,minimum width=4.6cm] (cp) at (12.0,-5.6) {CaseState checkpoints\\(SQLite, after every node)};

% arrows main flow
\draw[arr] (in) -- node[lab,above=2pt]{text} (ext);
\draw[arr] (ext) -- node[lab,above=2pt]{fields} (plan);
\draw[arr] (plan) -- node[lab,above=2pt]{plan} (ret);
\draw[arr] (ret) -- node[lab,above=2pt]{evidence} (ver);
\draw[arr] (ver) -- node[lab,right=2pt]{checks\{ \}} (gr);
\draw[arr] (gr) -- node[lab,above=2pt]{grade} (dec);
\draw[arr] (dec) -- node[lab,above=2pt]{yes / limit} (gd);
\draw[arr] (gd) -- node[lab,above=2pt]{verdict} (out);
\draw[arr,dashed,BrickRed] (dec.north) -- (dec.north |- 0,-1.45) -| node[lab,pos=0.3,below=2pt]{no: re-plan (iteration $<3$)} (plan.south);
% tools
\draw[arr,<->] (ret.north) -- node[lab,left=2pt]{call / result + source ID} (ret.north |- tools.south);
\draw[arr] (tools.north -| reg) -- (reg.south);
\draw[arr] (tools.north -| docs) -- (docs.south);
\draw[arr] (tools.north -| cen) -- (cen.south);
% HITL
\draw[arr] (gd) -- node[lab,right=2pt,align=left]{\texttt{escalate\_to\_human}:\\insufficient, uncited\\or grade $<0.6$} (q);
\draw[arr] (q) -- node[lab,above=2pt]{case + reason} (hum);
\draw[arr] (hum) -- node[lab,left=2pt]{final decision + note} (out);
% checkpoints
\draw[arr,dotted,gray] (gr.south) -- node[lab,right=2pt]{state} (gr.south |- cp.north);

% legend
\node[slm,minimum width=1.2cm,minimum height=0.5cm,font=\scriptsize] (l1) at (0.2,5.0) {SLM step};
\node[agent,minimum width=1.2cm,minimum height=0.5cm,font=\scriptsize,right=0.25cm of l1] (l2) {Agent};
\node[rule,minimum width=1.2cm,minimum height=0.5cm,font=\scriptsize,below=0.25cm of l1] (l3) {Rules};
\node[human,minimum width=1.2cm,minimum height=0.5cm,font=\scriptsize,right=0.25cm of l3] (l4) {Human};
\end{tikzpicture}}
\caption{Architecture of LoomTrace: six-node state graph, tools, data stores and human review}
\label{fig:arch}
\end{figure}

\section{Component Description}

Each component in Figure~\ref{fig:arch} is described below in terms of its responsibility, its inputs
and outputs, the method it uses, and its behaviour on failure.

\subsection{Product Claim Input and Extract Node}

The input is the text of a product label or seller claim, optionally accompanied by a label image. The
\emph{Extract} node calls the \texttt{extract\_label} tool, which produces a typed \texttt{LabelFields}
object with the fields \texttt{product}, \texttt{region}, \texttt{state}, \texttt{material},
\texttt{weaver\_id}, \texttt{batch\_id} and a Boolean \texttt{gi\_claim}. If an image is provided and the
EasyOCR package~\cite{easyocr} is installed, the image text is recognised first. If an API key is
configured, a small language model (SLM), Llama~3.1~8B~\cite{llama3}, is asked to return the fields as
JSON. Otherwise, and whenever the model call fails, a deterministic parser reads ``key: value'' lines,
recognises identifier patterns such as \texttt{WV-KAN-001} and \texttt{B-KAN-0010}, and falls back to
searching the text for known product names. Missing fields are left empty rather than guessed, so that
later nodes can treat them as unknown.

\subsection{Planner}

The planner decides which checks the claim requires and what to search for. Three checks are always
planned (product registered, GI document support and census plausibility); region, state and material
checks are planned only when the corresponding field is present; and four producer checks (weaver
registered, weaver in the correct cluster, batch found, batch woven by the claimed weaver) are always
planned so that their absence is visible. On the first iteration the planner requests an exact match in
the registry and an unfiltered document search using the product name and region as the query. On a
re-plan it resolves the product name to its registered form by matching the longest known alias, and
builds a sharper query from the canonical name, region, state, registered materials and words that
occur in statements of case (``geographical area weavers woven''), filtered to that product. If the
product cannot be resolved, the second iteration records the product check as failed rather than
looping indefinitely.

\subsection{Retriever Agent and Tools}

The retriever agent executes the plan by calling the tools listed in Table~\ref{tab:tools} in
Chapter~6: \texttt{gi\_registry\_lookup}, \texttt{gi\_doc\_search}, \texttt{census\_lookup} and, when the
identifiers are present, \texttt{weaver\_lookup} and \texttt{provenance\_lookup}. Every tool returns a
typed \texttt{ToolResult} containing a success flag, the data, a source identifier, a note and the
elapsed time. The retriever appends each result to the evidence list in the state, tagged with the
iteration number. Tool failures, such as an unknown weaver identifier, are returned as results with
\texttt{ok=false} and an explanatory note; they are never raised as exceptions, so that the verifier can
reason about them.

\subsection{Verifier Agent}

The verifier agent judges the evidence of the current iteration. It runs up to ten rule-based checks
and assigns each a status of \emph{pass}, \emph{fail} or \emph{unknown}, together with a plain-language
reason and the source identifier on which the status rests. For example, the region check passes if
the claimed region matches one of the registered districts, fails if it does not, and is unknown if
the product could not be resolved; for GIs registered for a whole state, such as Muga Silk, the check
compares the state instead. A claimed weaver who is not in the registry yields \emph{unknown} for the
batch--weaver match rather than \emph{fail}, because a mismatch cannot be established against an
unverified identity. Deterministic rules are used for these checks because they are exact and
repeatable, which a small model is not.

\subsection{Grader}

The grader scores the evidence before any verdict is produced. It computes completeness $C$, the
fraction of checks that are not unknown, and the citation fraction $Q$, the fraction of determinate
checks that carry a source identifier. If no check has failed, the grade is $g = 0.7C + 0.3Q$; if a
check has failed, the grade is $g = 0.5 + 0.5Q$, because a cited contradiction is strong evidence even
when other checks are incomplete. The grader then decides whether to loop: it requests a re-plan if the
product check or the document-support check is unknown, a product name is present, and fewer than three
iterations have been used. Otherwise it passes control to the guardrail.

\subsection{Guardrail and Human Review Gate}

The guardrail converts the checks into a verdict: \emph{Conflicting evidence} if any check failed,
\emph{Insufficient evidence} if any check is unknown, and \emph{Supported} only if every check passed.
It then enforces two conditions: a decided verdict must have at least one citation, and its grade must
be at least 0.6; if either condition fails, the verdict is downgraded to \emph{Insufficient evidence}.
Every \emph{Insufficient evidence} verdict is passed to the \texttt{escalate\_to\_human} tool, which
inserts the case, the extracted claim and the reason (for example, ``missing: weaver\_registered,
batch\_found'') into the review queue. The guardrail also writes the case to the case history.

\subsection{Review Queue and Human Reviewer}

The review queue is a table in the SQLite database. A reviewer sees each open case with its details
and reason, writes a note, and records one of three decisions (confirm authentic, not authentic, cannot
confirm). The decision is stored in a separate table and becomes the final verdict for the case, and
the requester can see it, with the note, by entering the case reference. A requester can also send a
case that was decided automatically to the queue if they disagree with the result.

\subsection{Data Stores and Checkpoints}

The four data stores are the GI registry (fifteen products), the GI passage index (21 passages), the
census table (31 rows) and the simulated weaver and batch records (180 and 360 records), described in
Section~3.4. The \texttt{CaseState} checkpoint store is an SQLite table to which the full state is
written after every node, giving each case its own memory and a durable history.

\section{Selected Agentic Patterns and Justification}

Table~\ref{tab:patterns} lists the patterns used, where each is applied, why it was chosen and which
alternative was rejected.

\begin{table}[H]
\centering
\caption{Agentic patterns and design justification}
\label{tab:patterns}
\small
\begin{tabularx}{\linewidth}{@{}>{\raggedright\arraybackslash}p{2.4cm} >{\raggedright\arraybackslash}p{2.4cm} >{\raggedright\arraybackslash}X >{\raggedright\arraybackslash}p{3.0cm}@{}}
\toprule
\textbf{Pattern} & \textbf{Where} & \textbf{Reason} & \textbf{Alternative rejected} \\ \midrule
Tool use & Retriever agent & Facts must come from registry, census and producer records, each with a source identifier & Answering from model memory \\
Agentic (corrective) RAG & Planner, retriever, grader & Labels use aliases; a weak first search must be detected and repeated with a sharper, filtered query~\cite{yan2024crag} & Single-shot retrieval \\
Plan-and-execute with re-planning & Planner, grader loop & Applicable checks depend on which fields are present; steps are known, so they can be bounded and traced~\cite{wang2023plansolve} & Free-form ReAct loop~\cite{yao2023react} \\
Supervisor multi-agent & Graph routes between retriever and verifier & One narrow responsibility per node: gathering evidence is separated from judging it & One node doing both \\
Reflection (evidence grading) & Grader & Verdicts must not be issued on incomplete evidence~\cite{shinn2023reflexion} & No self-check \\
HITL and guardrails & Guardrail, review queue & A wrong verdict harms a producer or a buyer; abstention is safer than a guess & Fully automatic decision \\
\bottomrule
\end{tabularx}
\end{table}

The main design decision was to use an explicit graph with deterministic checks rather than a
free-form ReAct agent. In a free-form agent, the model chooses which tool to call next and when to stop,
which makes the number of steps unbounded and the trace harder to audit. In LoomTrace the sequence of
nodes is fixed, only the plan contents and the loop decision vary, and the loop has a hard limit. A
second decision was to keep model calls narrow. Small models are fast and inexpensive but are weaker at
long, open-ended reasoning~\cite{liu2024lost}; giving the model only extraction and explanation, while
rules perform the factual checks, keeps the system correct even when no model is available.

\section{Control Flow and Agent State}

Control moves through the graph as follows. The graph starts at \emph{Extract}, then runs
\emph{Planner}, \emph{Retriever}, \emph{Verifier} and \emph{Grader} in sequence. A conditional edge
after the grader routes to the planner when the grader's decision is ``re-plan'' and to the guardrail
when it is ``finish''; the guardrail leads to the end of the graph. The loop terminates under exactly
three conditions: (i) the product and document checks are both determinate; (ii) no product name was
extracted, so re-planning cannot help; or (iii) the iteration counter has reached the limit of three
iterations. In addition, LangGraph's recursion limit is set to 40 node executions as a second,
independent bound. A full run without re-planning executes 6 nodes, and each re-plan adds 4 nodes.

Table~\ref{tab:state} lists the fields of the shared state.

\begin{table}[H]
\centering
\caption{Fields of the shared CaseState}
\label{tab:state}
\small
\begin{tabularx}{\linewidth}{@{}>{\ttfamily}l >{\raggedright\arraybackslash}X >{\raggedright\arraybackslash}p{2.6cm}@{}}
\toprule
\normalfont\textbf{Field} & \textbf{Content} & \textbf{Written by} \\ \midrule
case\_id & Unique reference, e.g.\ \texttt{CASE-94767A68} & Input \\
label\_text, image & Raw claim text and optional image bytes & Input \\
fields & Extracted \texttt{LabelFields} & Extract \\
plan, query, resolve\_aliases & Checks to run, search query, alias-resolution flag & Planner \\
canonical & Registered product name, once resolved & Retriever \\
evidence[ ] & Tool results with source IDs and iteration (append-only) & All tool-calling nodes \\
checks\{ \} & Status, reason and source per check & Verifier \\
grade, grade\_detail & Score 0--1, completeness, consistency, decision & Grader \\
iteration & Planning iterations used (0--2) & Extract, Grader \\
verdict, confidence, citations, rationale & Final output & Guardrail \\
escalated & Review-queue entry, if any & Guardrail \\
trace[ ] & One entry per node with a summary message (append-only) & All nodes \\
\bottomrule
\end{tabularx}
\end{table}

The \texttt{evidence} and \texttt{trace} fields are declared with an additive reducer, so that each
node appends to them instead of replacing them. This preserves the evidence of every iteration and
makes the full history of a case available to the verifier, the interface and the checkpoint store.

\section{Data Flow}

The following example traces one real input through the system. The label reads ``Product: Kanjivaram
Silk Saree; Region: Thiruvallur; State: Tamil Nadu; Material: mulberry silk; Weaver ID: WV-KAN-001;
Batch ID: B-KAN-0010''.

\begin{enumerate}[itemsep=2pt]
\item \textbf{Extract} produces the six fields; the product is the alias ``Kanjivaram Silk Saree''.
\item \textbf{Planner (iteration 0)} plans ten checks and the query ``Kanjivaram Silk Saree
Thiruvallur'', with exact registry matching.
\item \textbf{Retriever} calls five tools. The registry lookup fails (no exact match); the document
search returns five passages, two from the Kancheepuram corpus and three from other products; the
census, weaver and batch lookups succeed.
\item \textbf{Verifier} reports 5 pass, 0 fail and 5 unknown: the product, region, state, material and
cluster checks cannot be made without a registry entry.
\item \textbf{Grader} computes completeness 0.5 and citation fraction 1.0, so $g = 0.65$; because the
product check is unknown, it requests a re-plan.
\item \textbf{Planner (iteration 1)} resolves the alias to ``Kancheepuram Silk'' and issues a query
filtered to that product.
\item \textbf{Retriever} now finds registry entry GI~15 and three Kancheepuram passages.
\item \textbf{Verifier} reports 10 pass, 0 fail, 0 unknown; \textbf{Grader} computes $g = 1.0$ and
finishes.
\item \textbf{Guardrail} issues \emph{Supported} with seven citations: three GI passages, the registry
entry, the Tamil Nadu census row, and the weaver and batch records.
\end{enumerate}

The case executed ten nodes, made eleven tool calls and wrote ten checkpoints. The complete trace of
this run is given in Listing~\ref{lst:replan} in Chapter~5.
